\documentclass[11pt,a4paper]{article}
\usepackage[T1]{fontenc}
\usepackage[utf8]{inputenc}
\usepackage{lmodern,microtype,amsmath,amssymb,amsthm,mathtools,booktabs}
\usepackage[a4paper,margin=24mm,headheight=17pt]{geometry}
\usepackage{hyperref,xcolor,fancyhdr,graphicx}
\definecolor{accent}{HTML}{8F473D}
\hypersetup{colorlinks=true,linkcolor=accent,urlcolor=accent,pdftitle={Quantum Mechanics in Three Dimensions: Chapter 4 Worked Solutions}}
\pagestyle{fancy}\fancyhf{}\lhead{\small Quantum Mechanics / Chapter 4}
\rhead{\small Kumar Azad}\cfoot{\thepage}
\numberwithin{equation}{section}
\newcommand{\dd}{\,\mathrm d}
\newcommand{\ii}{\mathrm i}
\newcommand{\ee}{\mathrm e}
\newcommand{\ket}[1]{\lvert #1\rangle}
\newcommand{\bra}[1]{\langle #1\rvert}
\newcommand{\avg}[1]{\langle #1\rangle}
\newcommand{\braket}[2]{\langle #1\mid #2\rangle}
\newcommand{\comm}[2]{[#1,#2]}
\newcommand{\problem}[2]{\subsection*{Problem #1: #2}\addcontentsline{toc}{subsection}{Problem #1: #2}}
\newcommand{\answer}[1]{\par\smallskip\noindent\textcolor{accent}{\textbf{Result.}} #1\par\smallskip}
\newcommand{\note}[1]{\par\smallskip\noindent\textit{Check / scope.} #1\par\smallskip}
\newcommand{\aB}{a_{\mathrm B}}
\newcommand{\sh}{\hbar}
\setlength{\parindent}{0pt}\setlength{\parskip}{5pt}
\title{\textcolor{accent}{\large MATHEMATICAL PHYSICS · COURSEWORK COMPANION}\\[9pt]
\textbf{\Huge Quantum Mechanics\\in Three Dimensions}\\[11pt]
{\large Independent worked solutions to Chapter 4}\\[5pt]
{\normalsize Problems 4.1--4.76}}
\author{Kumar Azad\\\small\url{https://krazad0.github.io/QM-CH1-SOLUTION/chapter4/}}
\date{November 2026}
\begin{document}
\maketitle
\begin{abstract}
This educational manuscript develops independent stepwise solutions to the
76 numbered exercises in Chapter 4 of Griffiths and Schroeter,
\emph{Introduction to Quantum Mechanics}, third edition.
The topics include the three-dimensional Schrödinger equation,
spherical Bessel functions and spherical harmonics,
hydrogenic Coulomb spectra, orbital and spin angular momentum,
addition of spins, electromagnetic minimal coupling,
gauge covariance, and the Aharonov--Bohm effect.
Original descriptions paraphrase the textbook: consult the book for
the authoritative problem wording, numerical data, and figures.
The academic-style presentation is \textbf{not an arXiv publication
or a substitute for the textbook}.
\end{abstract}
\tableofcontents

\section{Notation and physical framework}
For a particle of mass $m$ in three dimensions,
\begin{equation}
\hat H=-\frac{\sh^2}{2m}\nabla^2+V(\mathbf r),\qquad
\hat{\mathbf p}=-\ii\sh\nabla,\qquad
[\hat r_i,\hat p_j]=\ii\sh\delta_{ij}.
\end{equation}
For a central potential, write
$\psi_{n\ell m}(\mathbf r)=R_{n\ell}(r)Y_\ell^m(\theta,\phi)$,
with
\begin{equation}
\int_0^\infty |R_{n\ell}(r)|^2r^2\dd r=1,\quad
\int|Y_\ell^m|^2\dd\Omega=1,\quad
\hat L^2Y_\ell^m=\sh^2\ell(\ell+1)Y_\ell^m.
\end{equation}
Use the Condon--Shortley phase convention for $Y_\ell^m$.
For hydrogen, set $\aB=4\pi\epsilon_0\sh^2/(\mu e^2)$
and $E_n=-\sh^2/(2\mu\aB^2n^2)$ when reduced mass
$\mu$ is used. Where the textbook neglects nuclear recoil,
set $\mu\simeq m_e$. For spin, $S_i=\sh\sigma_i/2$.
All numerical eigenvalue results specify units and numerical cutoffs;
a discretized spherical wall is not ordinary unconfined hydrogen.

\section{Three dimensions and spherical radial equations (4.1--4.11)}
\problem{4.1}{Canonical commutators and three-dimensional Ehrenfest theorem}
For smooth test functions $f(\mathbf r)$,
\begin{align}
[r_i,p_j]f&=-\ii\sh r_i\partial_j f+
\ii\sh\partial_j(r_if)=\ii\sh\delta_{ij}f,\\
[r_i,r_j]&=0,&[p_i,p_j]&=0.
\end{align}
Thus $[r_i,p_j]=\ii\sh\delta_{ij}$, including
$[x,p_y]=0$ and $[x,p_x]=\ii\sh$.
For $H=\sum_jp_j^2/(2m)+V$,
$[H,r_i]=-\ii\sh p_i/m$ and
$[H,p_i]=\ii\sh\partial_iV$. Apply
$\dd\langle Q\rangle/\dd t=(\ii/\sh)
\langle[H,Q]\rangle$:
\begin{equation}
\boxed{\frac{\dd\langle\mathbf r\rangle}{\dd t}
=\frac{\langle\mathbf p\rangle}{m},
\qquad
\frac{\dd\langle\mathbf p\rangle}{\dd t}
=-\langle\nabla V\rangle}.
\end{equation}

\problem{4.2}{Infinite cubic well and accidental degeneracy}
Inside $0<x,y,z<a$, separate
$\psi=X(x)Y(y)Z(z)$, each factor satisfying
a one-dimensional infinite-well equation.
The normalized stationary states and energies are
\begin{align}
\boxed{\psi_{n_xn_yn_z}(\mathbf r)}
&=\boxed{\left(\frac2a\right)^{3/2}
\sin\frac{n_x\pi x}{a}
\sin\frac{n_y\pi y}{a}
\sin\frac{n_z\pi z}{a}},\\
\boxed{E_{n_xn_yn_z}}&=
\frac{\pi^2\sh^2}{2ma^2}(n_x^2+n_y^2+n_z^2).
\end{align}
Writing $E_*=\pi^2\sh^2/(2ma^2)$, distinct
levels begin
\begin{equation}
\begin{array}{c|rrrrrr}
\text{level}&E_1&E_2&E_3&E_4&E_5&E_6\\\hline
E/E_*&3&6&9&11&12&14\\
\text{degeneracy}&1&3&3&3&1&6
\end{array}.
\end{equation}
The fourteenth distinct level is
$\boxed{E_{14}=27E_*}$ with
$\boxed{4}$ states: three permutations of $(1,1,5)$
and the distinct state $(3,3,3)$.
The latter is an accidental degeneracy,
not explained by permuting the three axes.

\problem{4.3}{Reconstructing central potentials from wavefunctions}
For a real radial wavefunction $\psi(r)$,
$\nabla^2\psi=\psi''+(2/r)\psi'$, so the
stationary equation gives
\begin{equation}
\boxed{V(r)=E+\frac{\sh^2}{2m}
\frac{\psi''+2\psi'/r}{\psi}}.
\end{equation}
(a) If $\psi=A\ee^{-r/a}$,
$\psi'/\psi=-1/a$,
$\psi''/\psi=1/a^2$, and
\begin{equation}
V(r)=E+\frac{\sh^2}{2m}
\left(\frac1{a^2}-\frac2{ar}\right).
\end{equation}
As $V(\infty)=0$,
$\boxed{E=-\sh^2/(2ma^2)}$ and
$\boxed{V=-\sh^2/(mar)}$.
(b) If $\psi=A\ee^{-r^2/a^2}$,
$\psi'/\psi=-2r/a^2$,
$\psi''/\psi=-2/a^2+4r^2/a^4$, yielding
$\nabla^2\psi/\psi=-6/a^2+4r^2/a^4$.
Setting $V(0)=0$ gives
\begin{equation}
\boxed{E=3\sh^2/(ma^2),\qquad
V(r)=2\sh^2r^2/(ma^4)}.
\end{equation}
This is an isotropic oscillator of frequency
$\omega=2\sh/(ma^2)$.

\problem{4.4}{Constructing and normalizing spherical harmonics}
With the Condon--Shortley convention,
$P_0=1$, $P_2^1(x)=-3x\sqrt{1-x^2}$ and
\begin{equation}
\boxed{Y_0^0=\frac1{\sqrt{4\pi}},
\qquad Y_2^1=-\sqrt{\frac{15}{8\pi}}\sin\theta
\cos\theta\,\ee^{\ii\phi}}.
\end{equation}
For $Y_{00}$,
$\int|Y_{00}|^2\dd\Omega=4\pi/(4\pi)=1$.
For $Y_{21}$,
$\int_0^{2\pi}\dd\phi=2\pi$ and
$\int_0^\pi\sin^3\theta\cos^2\theta\dd\theta
=4/15$, hence its norm is
$(15/(8\pi))(2\pi)(4/15)=1$.
Orthogonality follows by integrating the
factor $\ee^{\ii\phi}$ over $\phi$.

\problem{4.5}{Why the second angular solution is unacceptable}
For $\ell=m=0$ the polar angular equation is
$\frac1{\sin\theta}\frac{\dd}{\dd\theta}
(\sin\theta\,\Theta')=0$. Integrating gives
$\sin\theta\,\Theta'=A$, so
\begin{equation}
\Theta(\theta)=A\ln\tan\frac\theta2+B.
\end{equation}
Although the logarithm is integrable with
the $\sin\theta$ measure, it diverges at
$\theta=0,\pi$, and is not a regular smooth
angular eigenfunction on the sphere.
It is excluded by the regularity domain of
the self-adjoint angular Laplacian; the
acceptable $\ell=0$ solution is constant.
\note{Divergence at a coordinate pole,
not failure of square integrability alone,
is the key domain condition.}

\problem{4.6}{Conjugation identity of spherical harmonics}
Using the real associated Legendre function
and standard normalization,
$Y_\ell^m=N_{\ell m}P_\ell^m(\cos\theta)\ee^{\ii m\phi}$.
For integer $m>0$,
\begin{equation}
P_\ell^{-m}(x)=(-1)^m
\frac{(\ell-m)!}{(\ell+m)!}P_\ell^m(x),\quad
N_{\ell,-m}=
\sqrt{\frac{2\ell+1}{4\pi}
\frac{(\ell+m)!}{(\ell-m)!}}.
\end{equation}
Multiplying these and reversing the
azimuthal exponential gives
\begin{equation}
\boxed{Y_\ell^{-m}=(-1)^m(Y_\ell^{m})^*}.
\end{equation}

\problem{4.7}{Highest-weight harmonics and the explicit $Y_3^2$}
Rodrigues' expression gives
$P_\ell^\ell(x)=(-1)^\ell(2\ell-1)!!
(1-x^2)^{\ell/2}$.
Therefore
\begin{equation}
\boxed{Y_\ell^\ell=
(-1)^\ell
\sqrt{\frac{2\ell+1}{4\pi}
\frac{(2\ell-1)!!^2}{(2\ell)!}}\,
\sin^\ell\theta\,\ee^{\ii\ell\phi}}.
\end{equation}
For $\ell=3$, $P_3^2(x)=15x(1-x^2)$,
and
\begin{equation}
\boxed{Y_3^2=
\sqrt{\frac{105}{32\pi}}\,
\sin^2\theta\cos\theta\,\ee^{2\ii\phi}}.
\end{equation}
Substitution into the angular Laplacian
$\Delta_\Omega=
(\sin\theta)^{-1}\partial_\theta
(\sin\theta\,\partial_\theta)
+(\sin^2\theta)^{-1}\partial_\phi^2$
gives $\Delta_\Omega Y_\ell^m
=-\ell(\ell+1)Y_\ell^m$;
in the second case the eigenvalue is $-12$.

\problem{4.8}{Legendre orthogonality from Rodrigues' formula}
The Legendre functions obey
$\frac{\dd}{\dd x}[(1-x^2)P_\ell']
+\ell(\ell+1)P_\ell=0$.
Multiply the $\ell$ equation by $P_j$,
the $j$ equation by $P_\ell$, subtract,
and integrate over $[-1,1]$:
\begin{equation}
[\ell(\ell+1)-j(j+1)]
\int_{-1}^1P_\ell P_j\dd x
=\left[(1-x^2)(P_jP_\ell'-P_\ell P_j')\right]_{-1}^1=0.
\end{equation}
For $\ell\ne j$ the overlap vanishes.
To normalize, Rodrigues' formula gives
$P_\ell=(2^\ell\ell!)^{-1}\frac{\dd^\ell}{\dd x^\ell}
(x^2-1)^\ell$. Integrating by parts $\ell$
times, boundary terms vanish, leaving
\begin{equation}
\int_{-1}^1P_\ell^2\dd x
=\frac{(2\ell)!}{2^{2\ell}(\ell!)^2}
\int_{-1}^1(1-x^2)^\ell\dd x
=\frac{2}{2\ell+1}.
\end{equation}
Hence
$\boxed{\int_{-1}^1P_\ell P_j\dd x=
2\delta_{\ell j}/(2\ell+1)}$.

\problem{4.9}{Spherical Neumann functions near the origin}
The irregular spherical Bessel functions
are $n_0(x)=-\cos x/x$, $n_1(x)=-\cos x/x^2-\sin x/x$
and $n_{\ell+1}=(2\ell+1)n_\ell/x-n_{\ell-1}$.
Therefore
\begin{equation}
\boxed{n_1(x)=-\frac{\cos x}{x^2}
-\frac{\sin x}{x},\quad
n_2(x)=
-\left(\frac3{x^3}-\frac1x\right)\cos x
-\frac{3\sin x}{x^2}}.
\end{equation}
With $\sin x=x-x^3/6+\cdots$ and
$\cos x=1-x^2/2+\cdots$,
\begin{equation}
\boxed{n_1(x)=-x^{-2}-\tfrac12+O(x^2),\qquad
n_2(x)=-3x^{-3}-\tfrac1{2x}+O(x)}.
\end{equation}
The irregular solutions diverge at $r=0$;
for a finite central potential and
physical regularity one selects $j_\ell$.

\problem{4.10}{Spherical hard-wall energies at $\ell=1$}
For free radial motion $R(r)=Aj_1(kr)$,
$j_1(z)=\sin z/z^2-\cos z/z$.
It obeys $R''+(2/r)R'+[k^2-2/r^2]R=0$
by substitution (the $\ell=1$ radial equation).
A hard wall at $r=a$ imposes $j_1(ka)=0$:
\begin{equation}
\sin z-z\cos z=0
\quad\Longleftrightarrow\quad
\boxed{\tan z=z},\qquad z=ka>0.
\end{equation}
The first three positive roots are
$z_1\simeq4.493409$, $z_2\simeq7.725252$,
$z_3\simeq10.904122$, giving
$\boxed{E_{N1}=\sh^2 z_N^2/(2ma^2)}$.
For large $N$, $z_N\sim(N+1/2)\pi$ and
\begin{equation}
\boxed{E_{N1}\sim\frac{\pi^2\sh^2}{2ma^2}
(N+\tfrac12)^2}.
\end{equation}
The $\ell=1$ magnetic degeneracy is three.

\problem{4.11}{Finite spherical well and the threshold for binding}
For $V=-V_0$ inside $r<a$ and $V=0$
outside, write $u=rR$, $E=-\sh^2\kappa^2/(2m)$,
$k^2=2m(V_0+E)/\sh^2$. For $\ell=0$,
regularity at the origin gives
$u_{\rm in}=A\sin kr$, while
$u_{\rm out}=B\ee^{-\kappa r}$.
Matching $u,u'$ at $a$ gives
\begin{equation}
\boxed{k\cot(ka)=-\kappa},\qquad
k^2+\kappa^2=\frac{2mV_0}{\sh^2}.
\end{equation}
At the first threshold $\kappa\to0^+$,
$ka\to\pi/2$, so a bound level exists
only when $ka$ can reach at least $\pi/2$:
\begin{equation}
\boxed{V_0a^2>\frac{\pi^2\sh^2}{8m}}.
\end{equation}
At equality the zero-energy threshold
state is not square-integrable outside;
there is no normalizable bound state.
This contrasts with one-dimensional
finite wells, which bind at arbitrarily
small integrated attraction.

\section{Hydrogen atom and Coulomb observables (4.12--4.20)}
\problem{4.12}{Unnormalized hydrogen radial functions at $n=3$}
With $\rho=2r/(3\aB)$,
the hydrogen radial equation admits
$R_{n\ell}=C\,\rho^\ell\ee^{-\rho/2}
L_{n-\ell-1}^{2\ell+1}(\rho)$.
For $n=3$ the relevant associated Laguerre
polynomials are $L_2^1(\rho)=
\frac12(\rho^2-6\rho+6)$,
$L_1^3(\rho)=4-\rho$, and $L_0^5=1$.
Therefore, up to overall normalization,
\begin{align}
\boxed{R_{30}}&\propto
(6-6\rho+\rho^2)\ee^{-\rho/2},\\
\boxed{R_{31}}&\propto
\rho(4-\rho)\ee^{-\rho/2},\\
\boxed{R_{32}}&\propto
\rho^2\ee^{-\rho/2}.
\end{align}
The terminating polynomial degrees are respectively
$2,1,0$; all decay as $\ee^{-r/(3\aB)}$.

\problem{4.13}{Normalizing the $n=2$ hydrogen eigenfunctions}
For $\rho=r/\aB$, the radial functions take
the normalized forms
\begin{align}
\boxed{R_{20}(r)}&=\frac{1}{2\sqrt2\,\aB^{3/2}}
(2-\rho)\ee^{-\rho/2},\\
\boxed{R_{21}(r)}&=\frac1{2\sqrt6\,\aB^{3/2}}
\rho\ee^{-\rho/2}.
\end{align}
For example
$\int_0^\infty |R_{20}|^2r^2\dd r
=\frac18\int_0^\infty
\rho^2(2-\rho)^2\ee^{-\rho}\dd\rho
=\frac18(8-24+24)=1$.
For $R_{21}$, the integral equals
$\frac1{24}\int\rho^4\ee^{-\rho}\dd\rho=1$.
Combining with the normalized harmonics,
\begin{align}
\boxed{\psi_{200}}&=
\frac{(2-\rho)\ee^{-\rho/2}}{4\sqrt{2\pi}\,\aB^{3/2}},\\
\boxed{\psi_{210}}&=
\frac{\rho\cos\theta\,\ee^{-\rho/2}}
{4\sqrt{2\pi}\,\aB^{3/2}},\\
\boxed{\psi_{211}}&=
-\frac{\rho\sin\theta\,\ee^{\ii\phi}\ee^{-\rho/2}}
{8\sqrt{\pi}\,\aB^{3/2}},\\
\boxed{\psi_{21,-1}}&=
\frac{\rho\sin\theta\,\ee^{-\ii\phi}\ee^{-\rho/2}}
{8\sqrt{\pi}\,\aB^{3/2}}.
\end{align}
Angular orthonormality and the radial normalization
ensure unit three-dimensional norm.

\problem{4.14}{Laguerre polynomials and a higher radial solution}
With the ordinary Laguerre convention
$L_q(x)=\ee^x(q!)^{-1}
\dd^q(\ee^{-x}x^q)/\dd x^q$,
\begin{align}
L_0&=1,&L_1&=1-x,\\
L_2&=1-2x+\tfrac12x^2,&
L_3&=1-3x+\tfrac32x^2-\tfrac16x^3.
\end{align}
The radial polynomial for $n=5,\ell=2$ has
$n-\ell-1=2$ and parameter $2\ell+1=5$,
so
\begin{equation}
\boxed{L_2^5(\rho)=21-7\rho+\frac12\rho^2}.
\end{equation}
In the textbook's radial ansatz $u(\rho)=
\rho^{\ell+1}\ee^{-\rho/2}v(\rho)$,
the polynomial is
$\boxed{v(\rho)\propto21-7\rho+\rho^2/2}$.
Checking the series recurrence obtained by
substituting $v=\sum_jc_j\rho^j$ into
\begin{equation}
\rho v''+(2\ell+2-\rho)v'
+(n-\ell-1)v=0
\end{equation}
gives $(j+1)(j+2\ell+2)c_{j+1}=
(j+\ell+1-n)c_j$.
At $n=5,\ell=2$:
$c_1/c_0=-1/3$, $c_2/c_1=-1/14$,
$c_3=0$, confirming the polynomial.

\problem{4.15}{Hydrogen mean radii and Cartesian position moments}
The $1s$ density is
$|\psi_{100}|^2=\ee^{-2r/\aB}/(\pi\aB^3)$.
Using $\int_0^\infty r^j\ee^{-\beta r}\dd r
=j!/\beta^{j+1}$,
\begin{align}
\avg r&=\frac4{\aB^3}\int_0^\infty
r^3\ee^{-2r/\aB}\dd r
=\boxed{\tfrac32\aB},\\
\avg{r^2}&=\frac4{\aB^3}\int_0^\infty
r^4\ee^{-2r/\aB}\dd r
=\boxed{3\aB^2}.
\end{align}
Spherical symmetry implies
$\boxed{\avg x=0,\ \avg{x^2}=\aB^2}$.
For the $n=2,\ell=1,m=1$ state,
the standard radial integral gives
$\avg{r^2}=30\aB^2$.
Its angular density is proportional to
$\sin^2\theta$, for which
$\avg{\sin^2\theta\cos^2\phi}=2/5$;
hence
\begin{equation}
\boxed{\avg{x^2}_{211}=12\aB^2}.
\end{equation}

\problem{4.16}{Most probable hydrogenic ground-state radius}
The probability in a spherical shell of thickness
$\dd r$ is $P(r)\dd r$ with
\begin{equation}
P(r)=4\pi r^2|\psi_{100}(r)|^2
=\frac{4r^2}{\aB^3}\ee^{-2r/\aB}.
\end{equation}
Differentiate:
$P'(r)=\frac{8r}{\aB^3}\ee^{-2r/\aB}
(1-r/\aB)$.
The positive stationary point is
$r=\aB$; the density vanishes at $r=0$
and at infinity, so the maximum is
\begin{equation}
\boxed{r_{\rm most\ probable}=\aB}.
\end{equation}
The radial maximum is not the same
as the maximum of the spatial density,
which occurs at $r=0$.

\problem{4.17}{A double-commutator shortcut for $\langle zHz\rangle$}
Since $[H,z]=-\ii\sh p_z/m$ and
$[z,p_z]=\ii\sh$, one has
\begin{equation}
[z,[H,z]]=\frac{\sh^2}{m}I.
\end{equation}
Expanding the double commutator gives
$[z,[H,z]]=2zHz-z^2H-Hz^2$.
In the normalized ground eigenstate
$H\ket{100}=E_1\ket{100}$,
\begin{equation}
2\avg{zHz}-2E_1\avg{z^2}
=\frac{\sh^2}{m}.
\end{equation}
By Problem 4.15, $\avg{z^2}=\aB^2$ and
$E_1=-\sh^2/(2m\aB^2)$
(fixed-nucleus approximation), so
\begin{equation}
\boxed{\avg{zHz}=E_1\aB^2+\frac{\sh^2}{2m}=0}.
\end{equation}
No six-dimensional integral is necessary.

\problem{4.18}{Stationarity within the hydrogen $n=2$ eigenspace}
For the initial state
$\Psi(\mathbf r,0)=
(\psi_{211}+\psi_{21,-1})/\sqrt2$,
both components share energy
$E_2=-13.6/4\,{\rm eV}=-3.40\,{\rm eV}$.
Therefore
\begin{equation}
\boxed{\Psi(\mathbf r,t)=
\ee^{-\ii E_2t/\sh}\,
\frac{\psi_{211}(\mathbf r)+\psi_{21,-1}(\mathbf r)}
{\sqrt2}}.
\end{equation}
Using $Y_{11}\propto-\sin\theta\ee^{\ii\phi}$,
$Y_{1,-1}\propto+\sin\theta\ee^{-\ii\phi}$,
the angular part is proportional to
$-\ii\sin\theta\sin\phi$ (a $p_y$ orbital).
Only a global time-dependent phase appears:
$|\Psi(\mathbf r,t)|^2$ is stationary.
For Coulomb $V\propto1/r$, the virial theorem
gives $\avg V=2E_2$; thus
\begin{equation}
\boxed{\avg V=2E_2=-6.80\,{\rm eV}},
\end{equation}
independent of time.

\problem{4.19}{Hydrogenic scaling with nuclear charge $Z$}
For the fixed-nucleus Coulomb potential
$V_Z=-Ze^2/(4\pi\epsilon_0r)$,
replace $e^2$ by $Ze^2$ in the hydrogen
radius and energy formulas. Hence
\begin{equation}
\boxed{a_Z=\frac{\aB}{Z},\quad
E_n(Z)=-\frac{13.6\,Z^2}{n^2}\,{\rm eV},\quad
|E_1(Z)|=13.6Z^2\,{\rm eV},\quad
R_Z=Z^2R_H}.
\end{equation}
With finite nuclear mass, additionally
replace electron mass by the corresponding
electron--nucleus reduced mass.
Lyman wavelengths scale as $1/Z^2$:
for $Z=2$, the series is at one quarter
of the hydrogen wavelengths (extreme UV);
for $Z=3$, one ninth (extreme UV/soft-X-ray
boundary for the shortest wavelengths).
The classification depends on the exact
wavelength threshold adopted.

\problem{4.20}{Gravitational hydrogen analogy for the Earth and Sun}
For Earth mass $m_E$, solar mass $M$, and
reduced mass $\mu=m_EM/(m_E+M)$,
\begin{equation}
\boxed{V(r)=-\frac{GMm_E}{r},\quad
a_g=\frac{\sh^2}{\mu GMm_E},\quad
E_n=-\frac{\mu(GMm_E)^2}{2\sh^2n^2}}.
\end{equation}
Using $M\simeq1.9885\times10^{30}$ kg,
$m_E\simeq5.9722\times10^{24}$ kg,
$\sh=1.05457\times10^{-34}$ J s, and
$G\simeq6.6743\times10^{-11}$ SI units,
\begin{equation}
\boxed{a_g\simeq2.35\times10^{-138}\ {\rm m}}.
\end{equation}
A classical circular orbit at radius $r_0$ has
$E_{\rm cl}=-GMm_E/(2r_0)$; comparing with $E_n$
gives $r_0=a_gn^2$ and
\begin{equation}
\boxed{n=\sqrt{r_0/a_g}\simeq
2.52\times10^{74}}
\end{equation}
for one astronomical unit.
For $n\gg1$, the next-level energy loss is
$\Delta E\sim2|E_n|/n$, giving
$\boxed{\Delta E\simeq2.10\times10^{-41}\ {\rm J}}$.
The corresponding quantum wavelength
$\lambda=hc/\Delta E$ is
\begin{equation}
\boxed{\lambda\simeq9.46\times10^{15}\ {\rm m}
\simeq1.00\ {\rm light\ year}}.
\end{equation}
The near-one-light-year number reflects
the numerical relationship of the
Earth--Sun orbital parameters and
fundamental constants; the toy model
is not a realistic description of
gravitational radiation from the orbit.

\section{Angular momentum and elementary spin (4.21--4.34)}
\problem{4.21}{Ladder normalization for orbital angular momentum}
With $L_\pm=L_x\pm\ii L_y$ and
$L_\pm^\dagger=L_\mp$,
\begin{equation}
L_\mp L_\pm=L^2-L_z^2\mp\sh L_z.
\end{equation}
If $\ket{\ell m}$ is normalized,
$\|L_+\ket{\ell m}\|^2=
\bra{\ell m}L_-L_+\ket{\ell m}
=\sh^2[\ell(\ell+1)-m(m+1)]$.
Therefore
\begin{align}
\boxed{L_+\ket{\ell m}}
&=\boxed{\sh\sqrt{(\ell-m)(\ell+m+1)}
\ket{\ell,m+1}},\\
\boxed{L_-\ket{\ell m}}
&=\boxed{\sh\sqrt{(\ell+m)(\ell-m+1)}
\ket{\ell,m-1}}.
\end{align}
The top and bottom coefficients vanish at $m=\pm\ell$,
respectively.

\problem{4.22}{Orbital commutators and rotational invariance}
Since $L_z=xp_y-yp_x$,
\begin{align}
[L_z,x]&=[xp_y,x]-[yp_x,x]=\ii\sh y,&
[L_z,y]&=-\ii\sh x,&[L_z,z]&=0,\\
[L_z,p_x]&=\ii\sh p_y,&
[L_z,p_y]&=-\ii\sh p_x,&[L_z,p_z]&=0.
\end{align}
Insert $L_x=yp_z-zp_y$ into
$[L_z,L_x]$ and use the product rule
to obtain $\boxed{[L_z,L_x]=\ii\sh L_y}$.
The same commutators imply
$[L_z,x^2+y^2+z^2]=0$ and
$[L_z,p_x^2+p_y^2+p_z^2]=0$;
hence
$\boxed{[L_z,r]=[L_z,p^2]=0}$.
For central $V(r)$, $[L_i,V(r)]=0$
and $[L_i,p^2]=0$ for all $i$,
so $\boxed{[H,\mathbf L]=0}$ and
$[H,L^2]=0$.

\problem{4.23}{Quantum torque and orbital-angular-momentum conservation}
For $L_i=\epsilon_{ijk}r_jp_k$,
apply the generalized Ehrenfest theorem:
\begin{align}
\frac{\dd}{\dd t}\avg{L_i}
&=\frac{\ii}{\sh}
\epsilon_{ijk}\avg{[H,r_jp_k]}\\
&=\epsilon_{ijk}\avg{(p_j/m)p_k
-r_j\partial_kV}\\
&=-\epsilon_{ijk}\avg{r_j\partial_kV}.
\end{align}
The term with $p_jp_k$ vanishes
because it is symmetric in $j,k$.
Thus
\begin{equation}
\boxed{\frac{\dd\avg{\mathbf L}}{\dd t}
=\avg{\mathbf r\times(-\nabla V)}}.
\end{equation}
For $V=V(r)$, $\nabla V=V'(r)\hat{\mathbf r}$,
so the torque vanishes identically and
$\boxed{\dd\avg{\mathbf L}/\dd t=0}$.

\problem{4.24}{Differential representation of $L^2$}
For a test function $f(\theta,\phi)$,
\begin{equation}
L_\pm f=\pm\sh\ee^{\pm\ii\phi}
(\partial_\theta\pm\ii\cot\theta\,\partial_\phi)f.
\end{equation}
Apply $L_+$ to $L_-f$, differentiating
both the exponential and $\cot\theta$:
\begin{equation}
L_+L_-=-\sh^2\left(
\partial_\theta^2+\cot\theta\,\partial_\theta+
\cot^2\theta\,\partial_\phi^2+
\ii\partial_\phi\right).
\end{equation}
Using the algebraic identity
$L^2=L_+L_-+L_z^2-\sh L_z$
and $L_z=-\ii\sh\partial_\phi$,
the first-derivative terms cancel.
The coefficient of $\partial_\phi^2$
combines as $1+\cot^2\theta=\csc^2\theta$:
\begin{equation}
\boxed{L^2=-\sh^2
\left[\frac1{\sin\theta}\partial_\theta
(\sin\theta\,\partial_\theta)+
\frac1{\sin^2\theta}\partial_\phi^2\right]}.
\end{equation}

\problem{4.25}{Highest-weight spherical harmonic}
By the ladder algebra,
$\boxed{L_+Y_\ell^\ell=0}$.
Let $Y_\ell^\ell(\theta,\phi)=
F(\theta)\ee^{\ii\ell\phi}$.
Then $L_+Y=0$ reduces to
$F'-\ell\cot\theta\,F=0$ and hence
$F(\theta)=C\sin^\ell\theta$.
Normalize:
\begin{equation}
1=2\pi|C|^2\int_0^\pi
\sin^{2\ell+1}\theta\dd\theta,\quad
\int_0^\pi\sin^{2\ell+1}\theta\dd\theta
=\frac{2^{2\ell+1}(\ell!)^2}{(2\ell+1)!}.
\end{equation}
Choosing the standard phase gives
\begin{equation}
\boxed{Y_\ell^\ell=
(-1)^\ell\frac{\sqrt{(2\ell+1)!}}
{2^\ell\ell!\sqrt{4\pi}}\,
\sin^\ell\theta\ee^{\ii\ell\phi}}.
\end{equation}
This agrees with the associated-Legendre formula.

\problem{4.26}{Raising $Y_2^1$ to obtain $Y_2^2$}
Start from
$Y_2^1=-\sqrt{15/(8\pi)}
\sin\theta\cos\theta\,\ee^{\ii\phi}$.
Apply
$L_+=\sh\ee^{\ii\phi}
(\partial_\theta+\ii\cot\theta\,\partial_\phi)$:
\begin{equation}
L_+Y_2^1=
\sh\sqrt{\frac{15}{8\pi}}
\sin^2\theta\,\ee^{2\ii\phi}.
\end{equation}
The normalization factor from Problem 4.21
is $\sh\sqrt{(2-1)(2+1+1)}=2\sh$.
Therefore
\begin{equation}
\boxed{Y_2^2=
\sqrt{\frac{15}{32\pi}}\,
\sin^2\theta\,\ee^{2\ii\phi}}.
\end{equation}

\problem{4.27}{Rigid rotor spectrum of a diatomic molecule}
The center-of-mass distances are
$r_1=m_2a/(m_1+m_2)$ and
$r_2=m_1a/(m_1+m_2)$, so
\begin{equation}
I=m_1r_1^2+m_2r_2^2
=\boxed{\frac{m_1m_2}{m_1+m_2}a^2}.
\end{equation}
For a rigid rotor $H=L^2/(2I)$, giving
\begin{equation}
\boxed{E_\ell=\frac{\sh^2\ell(\ell+1)}{2I}},
\quad\boxed{\psi_{\ell m}=Y_\ell^m(\theta,\phi)},
\quad d_\ell=2\ell+1.
\end{equation}
For electric-dipole rotational transitions
$\Delta\ell=\pm1$, the emission frequencies are
\begin{equation}
\boxed{\nu_j=\frac{E_j-E_{j-1}}h
=\frac{\sh j}{2\pi I}},\quad j=1,2,\dots.
\end{equation}
Adjacent lines are separated by
$\Delta\nu=\sh/(2\pi I)$.
Using the approximate separation in the
textbook CO figure, $\Delta\tilde\nu\simeq
3.84\ {\rm cm^{-1}}$,
$I=\sh/(2\pi c\Delta\tilde\nu)$.
The $^{12}$C--$^{16}$O reduced mass is
about $6.86$ atomic mass units, giving
\begin{equation}
\boxed{a=\sqrt{I/\mu}\simeq1.13\times10^{-10}\ {\rm m}}.
\end{equation}
\textit{The line spacing is read approximately from a figure;
the quoted bond length is correspondingly approximate.}

\problem{4.28}{Why a classical spinning electron cannot explain spin}
For a uniform classical solid sphere,
$I=(2/5)m_er_c^2$, where
$r_c=e^2/(4\pi\epsilon_0m_ec^2)$.
Equating $I\omega=\sh/2$ gives
\begin{equation}
v_{\rm equator}=\omega r_c
=\boxed{\frac{5\sh}{4m_er_c}}.
\end{equation}
With $r_c\simeq2.818\times10^{-15}$ m,
\begin{equation}
\boxed{v_{\rm equator}\simeq5.1\times10^{10}
\ {\rm m/s}\simeq171\,c}.
\end{equation}
This impossible speed shows that electron
spin is intrinsic angular momentum, not the
classical rotation of a tiny extended sphere.

\problem{4.29}{Spin commutators and Pauli-product identity}
For spin one-half, $S_j=\sh\sigma_j/2$,
where
\begin{equation}
\sigma_x=\begin{pmatrix}0&1\\1&0\end{pmatrix},
\quad\sigma_y=\begin{pmatrix}0&-\ii\\\ii&0\end{pmatrix},
\quad\sigma_z=\begin{pmatrix}1&0\\0&-1\end{pmatrix}.
\end{equation}
Explicit multiplication gives
$\sigma_x\sigma_y=\ii\sigma_z$,
$\sigma_y\sigma_x=-\ii\sigma_z$,
and $\sigma_j^2=I$. Cyclic permutation yields
\begin{equation}
\boxed{\sigma_j\sigma_k=
\delta_{jk}I+\ii\sum_\ell
\epsilon_{jk\ell}\sigma_\ell},
\end{equation}
hence
$\boxed{[S_j,S_k]=\ii\sh\epsilon_{jk\ell}S_\ell}$.

\problem{4.30}{Measured spins and uncertainties for $(3i,4)$}
The state $\chi=A(3\ii,4)^{\mathsf T}$
has $A=1/5$ (real, positive).
With the Pauli matrices,
\begin{equation}
\boxed{\avg{S_x}=0,\quad
\avg{S_y}=-\frac{12\sh}{25},\quad
\avg{S_z}=-\frac{7\sh}{50}}.
\end{equation}
Since $S_j^2=\sh^2I/4$,
\begin{equation}
\boxed{\sigma_{S_x}=\frac\sh2,\quad
\sigma_{S_y}=\frac{7\sh}{50},\quad
\sigma_{S_z}=\frac{12\sh}{25}}.
\end{equation}
The products
$\sigma_{S_x}\sigma_{S_y}=7\sh^2/100$ and
$\sigma_{S_x}\sigma_{S_z}=6\sh^2/25$
equal $(\sh/2)|\avg{S_z}|$
and $(\sh/2)|\avg{S_y}|$, respectively.
The remaining inequality is automatic because
$\avg{S_x}=0$.

\problem{4.31}{Expectations of a general spin-one-half state}
For normalized $\chi=(a,b)^{\mathsf T}$,
$|a|^2+|b|^2=1$, matrix multiplication yields
\begin{equation}
\boxed{\avg{S_x}=\sh\,\mathrm{Re}(a^*b),\quad
\avg{S_y}=\sh\,\mathrm{Im}(a^*b),\quad
\avg{S_z}=\frac\sh2(|a|^2-|b|^2)}.
\end{equation}
Since all Pauli matrices square to $I$,
\begin{equation}
\boxed{\avg{S_x^2}=\avg{S_y^2}=\avg{S_z^2}
=\frac{\sh^2}{4}},\qquad
\boxed{\avg{S^2}=\frac{3\sh^2}{4}}.
\end{equation}
The spin standard deviations follow by
subtracting squared means from $\sh^2/4$.

\problem{4.32}{Eigenstates and probabilities for $S_y$}
Diagonalize
$S_y=(\sh/2)\begin{pmatrix}0&-\ii\\\ii&0\end{pmatrix}$:
\begin{equation}
\boxed{\lambda_\pm=\pm\sh/2,\qquad
\chi_{y,+}=\frac1{\sqrt2}\binom{1}{\ii},
\quad\chi_{y,-}=\frac1{\sqrt2}\binom{1}{-\ii}}.
\end{equation}
For a general normalized $\chi=(a,b)^T$,
\begin{equation}
\boxed{P_y(+)=\frac12|a-\ii b|^2,\quad
P_y(-)=\frac12|a+\ii b|^2}.
\end{equation}
Their sum is $|a|^2+|b|^2=1$.
Since $S_y^2=(\sh^2/4)I$, measuring
$S_y^2$ yields only $\boxed{\sh^2/4}$,
with probability $1$.

\problem{4.33}{Spin along an arbitrary spatial direction}
For unit $\hat{\mathbf n}=
(\sin\theta\cos\phi,\sin\theta\sin\phi,\cos\theta)$,
\begin{equation}
S_{\mathbf n}=\frac\sh2
\begin{pmatrix}
\cos\theta&\ee^{-\ii\phi}\sin\theta\\
\ee^{\ii\phi}\sin\theta&-\cos\theta
\end{pmatrix}.
\end{equation}
Because $(\boldsymbol\sigma\cdot\hat{\mathbf n})^2=I$,
the eigenvalues are $\pm\sh/2$.
Solving the first row using
$\cos\theta\mp1=\mp2\sin^2(\theta/2)$
gives normalized eigenspinors
\begin{equation}
\boxed{\chi_{\mathbf n,+}=
\binom{\cos(\theta/2)}
{\ee^{\ii\phi}\sin(\theta/2)}},
\quad
\boxed{\chi_{\mathbf n,-}=
\binom{\ee^{-\ii\phi}\sin(\theta/2)}
{-\cos(\theta/2)}}.
\end{equation}
They are orthogonal and have unit norm.

\problem{4.34}{Spin-one matrices}
Choose the ordered $S_z$ basis
$\ket{1,1}$, $\ket{1,0}$, $\ket{1,-1}$.
The ladder rule
$S_\pm\ket{1,m}=\sh
\sqrt{2-m(m\pm1)}\ket{1,m\pm1}$
gives
\begin{align}
\boxed{S_x}&=\frac\sh{\sqrt2}
\begin{pmatrix}0&1&0\\1&0&1\\0&1&0\end{pmatrix},\\
\boxed{S_y}&=\frac\sh{\sqrt2}
\begin{pmatrix}0&-\ii&0\\\ii&0&-\ii\\0&\ii&0\end{pmatrix},\\
\boxed{S_z}&=\sh\begin{pmatrix}
1&0&0\\0&0&0\\0&0&-1
\end{pmatrix}.
\end{align}
Direct products give
$[S_x,S_y]=\ii\sh S_z$,
and cyclic permutations, and
$S_x^2+S_y^2+S_z^2=2\sh^2I_3$.

\section{Spin dynamics, coupled spins, electromagnetic theory (4.35--4.45)}
\problem{4.35}{Spin measurement probabilities during Larmor precession}
Example 4.3 gives, for real initial coefficients,
$\chi(t)=(\cos(\alpha/2)\ee^{\ii\gamma B_0t/2},
\sin(\alpha/2)\ee^{-\ii\gamma B_0t/2})^T$.
The spin expectations are
\begin{equation}
\avg{S_x}=\frac\sh2\sin\alpha\cos\omega_Lt,\quad
\avg{S_y}=-\frac\sh2\sin\alpha\sin\omega_Lt,\quad
\avg{S_z}=\frac\sh2\cos\alpha,\quad \omega_L=\gamma B_0.
\end{equation}
For spin one-half,
$P(S_i=+\sh/2)=\frac12+\avg{S_i}/\sh$.
Therefore
\begin{align}
\boxed{P_x(+)}&=\frac12[1+\sin\alpha\cos(\omega_Lt)],\\
\boxed{P_y(+)}&=\frac12[1-\sin\alpha\sin(\omega_Lt)],\\
\boxed{P_z(+)}&=\frac12(1+\cos\alpha)=\cos^2(\alpha/2).
\end{align}

\problem{4.36}{Exact spin dynamics in an oscillating longitudinal field}
For $\mathbf B=B_0\cos(\omega t)\hat z$,
$H(t)=-\gamma B_0\cos(\omega t)S_z$,
so
\begin{equation}
\boxed{H(t)=-\frac{\gamma\sh B_0}{2}
\cos(\omega t)\begin{pmatrix}1&0\\0&-1\end{pmatrix}}.
\end{equation}
The Hamiltonian commutes with itself at all times,
so the evolution is an ordinary diagonal phase:
\begin{equation}
\boxed{U(t)=
\begin{pmatrix}\ee^{\ii\vartheta(t)}&0\\
0&\ee^{-\ii\vartheta(t)}\end{pmatrix}},
\quad
\vartheta(t)=\frac{\gamma B_0}{2\omega}\sin(\omega t).
\end{equation}
With initial $\chi_{x,+}=(1,1)^T/\sqrt2$,
$\chi(t)=
(\ee^{\ii\vartheta},\ee^{-\ii\vartheta})^T/\sqrt2$.
Projection onto $\chi_{x,-}=(1,-1)^T/\sqrt2$
gives amplitude $\ii\sin\vartheta$ and
\begin{equation}
\boxed{P(S_x=-\sh/2,t)=
\sin^2\left[\frac{\gamma B_0}{2\omega}\sin(\omega t)\right]}.
\end{equation}
Complete inversion requires the maximal accessible
$|\vartheta|\ge\pi/2$, so
$\boxed{|B_0|_{\min}=\pi|\omega|/|\gamma|}$,
for nonzero oscillation frequency.

\problem{4.37}{Confirming triplets and the singlet using ladder operators}
Write $\ket{1,0}=(\ket{\uparrow\downarrow}+
\ket{\downarrow\uparrow})/\sqrt2$.
With $S_-=S_-^{(1)}+S_-^{(2)}$,
\begin{equation}
S_-\ket{1,0}=
\frac1{\sqrt2}(\sh\ket{\downarrow\downarrow}
+\sh\ket{\downarrow\downarrow})
=\boxed{\sh\sqrt2\,\ket{1,-1}}.
\end{equation}
For the singlet
$\ket{0,0}=(\ket{\uparrow\downarrow}-
\ket{\downarrow\uparrow})/\sqrt2$,
the terms cancel,
$\boxed{S_\pm\ket{0,0}=0}$.
Since $S_z\ket{0,0}=0$ as well,
$S^2=S_+S_-+S_z^2-\sh S_z$
gives $\boxed{S^2\ket{0,0}=0}$.
For a triplet state the same expression gives
$S^2=2\sh^2$.

\problem{4.38}{Allowed baryon and meson spins}
For three spin-one-half constituents,
first couple two spins to $s_{12}=0$ or $1$.
Then $0\otimes\frac12=\frac12$,
while $1\otimes\frac12=\frac32\oplus\frac12$.
Thus
\begin{equation}
\boxed{\text{three quarks: }s=\frac12
\text{ or }\frac32}.
\end{equation}
For a quark and antiquark with spin one-half,
$\frac12\otimes\frac12=0\oplus1$:
\begin{equation}
\boxed{\text{two constituents: }s=0\text{ or }1}.
\end{equation}
These results assume zero orbital angular momentum.
Additional color, flavor and exchange constraints
are needed to classify physical hadrons.

\problem{4.39}{Clebsch--Gordan construction for two spin halves}
For $s_1=s_2=1/2$, the coupled triplet is
\begin{align}
\ket{1,1}&=\ket{\uparrow\uparrow},\\
\ket{1,0}&=\frac{\ket{\uparrow\downarrow}+
\ket{\downarrow\uparrow}}{\sqrt2},\\
\ket{1,-1}&=\ket{\downarrow\downarrow},
\end{align}
and the singlet is
$\ket{0,0}=(\ket{\uparrow\downarrow}-
\ket{\downarrow\uparrow})/\sqrt2$.
The corresponding Clebsch--Gordan coefficients
are $1,\;1/\sqrt2,\;1/\sqrt2,\;1$
for the triplet, and $+1/\sqrt2,-1/\sqrt2$
for the singlet. Each expansion is normalized.
Their inner products vanish across the
triplet/singlet sectors, confirming the
decomposition $2\otimes2=3\oplus1$.

\problem{4.40}{Probabilities for adding unequal spins}
(a) For $s_1=1,s_2=2$, the state $\ket{3,1}$
contains contributions with
$m_1+m_2=1$:
\begin{equation}
\ket{3,1}=
\sqrt{\frac25}\ket{1,1;2,0}
+\sqrt{\frac8{15}}\ket{1,0;2,1}
+\sqrt{\frac1{15}}\ket{1,-1;2,2}.
\end{equation}
Hence measuring $S_z^{(2)}$ yields
\begin{equation}
\boxed{m_2=0,1,2\quad\text{with probabilities }
2/5,\ 8/15,\ 1/15}.
\end{equation}
(b) A hydrogen electron in $\psi_{510}$
has orbital $\ell=1,m_\ell=0$; with spin
down $m_s=-1/2$. The possible total
$j$ values are $3/2$ and $1/2$.
The relevant Clebsch--Gordan coefficients
are $\sqrt{2/3}$ and $\sqrt{1/3}$,
respectively. Thus
\begin{equation}
\boxed{P(J^2=\tfrac{15}{4}\sh^2)=\frac23,\qquad
P(J^2=\tfrac34\sh^2)=\frac13}.
\end{equation}

\problem{4.41}{Total spin commutes with its $z$ component}
With $\mathbf S=\mathbf S_1+\mathbf S_2$,
$S^2=S_1^2+S_2^2+2\mathbf S_1\cdot\mathbf S_2$.
Each individual $S_a^2$ commutes with all
components of that particle's spin.
For the cross term,
\begin{equation}
[S_1\cdot S_2,S_{1k}]
=\sum_i[S_{1i},S_{1k}]S_{2i}
=\ii\sh\sum_{i,j}\epsilon_{ikj}S_{1j}S_{2i}.
\end{equation}
Rearranging indices yields the vector result
\begin{equation}
\boxed{[S^2,\mathbf S_1]
=2\ii\sh(\mathbf S_1\times\mathbf S_2)},
\quad
\boxed{[S^2,\mathbf S_2]
=-2\ii\sh(\mathbf S_1\times\mathbf S_2)}.
\end{equation}
Adding them gives $[S^2,\mathbf S]=0$,
and therefore $\boxed{[S^2,S_z]=0}$.
The corresponding commutator with an
\emph{individual} $S_{1z}$ is generally nonzero.

\problem{4.42}{Minimal coupling and the quantum Lorentz-force law}
Let $\boldsymbol\pi=\mathbf p-q\mathbf A$,
$\mathbf v=\boldsymbol\pi/m$ and
$H=\boldsymbol\pi^2/(2m)+q\phi$.
Using $[r_i,\pi_j]=\ii\sh\delta_{ij}$,
\begin{equation}
\boxed{\frac{\dd\avg{\mathbf r}}{\dd t}
=\frac{\avg{\boldsymbol\pi}}m=\avg{\mathbf v}}.
\end{equation}
The kinetic commutator is
$[\pi_i,\pi_j]=\ii q\sh\,\epsilon_{ijk}B_k$.
Applying Ehrenfest to $\pi_i$ and
including explicit $\partial_t\mathbf A$,
one obtains the symmetric-ordered result
\begin{equation}
m\frac{\dd\avg{\mathbf v}}{\dd t}
=q\avg{\mathbf E}
+\frac{q}{2m}
\avg{\boldsymbol\pi\times\mathbf B-
\mathbf B\times\boldsymbol\pi}.
\end{equation}
Writing $\boldsymbol\pi=\mathbf p-q\mathbf A$
recovers the equivalent textbook formula with
$\frac q{2m}(\mathbf p\times\mathbf B-
\mathbf B\times\mathbf p)-(q^2/m)\mathbf A\times\mathbf B$.
For spatially uniform $\mathbf E,\mathbf B$,
operator-ordering issues disappear:
\begin{equation}
\boxed{m\,\frac{\dd\avg{\mathbf v}}{\dd t}
=q(\mathbf E+\avg{\mathbf v}\times\mathbf B)}.
\end{equation}

\problem{4.43}{Landau levels plus confinement along the field}
With $\mathbf A=(B_0/2)(-y,x,0)$
and $\phi=Kz^2$, the electromagnetic fields are
\begin{equation}
\boxed{\mathbf B=B_0\hat z,\qquad
\mathbf E=-2Kz\hat z}.
\end{equation}
The Hamiltonian separates into a transverse
Landau problem and a $z$ oscillator:
\begin{equation}
H_\perp=\frac{(p_x+qB_0y/2)^2+
(p_y-qB_0x/2)^2}{2m},\qquad
H_z=\frac{p_z^2}{2m}+qKz^2.
\end{equation}
The transverse kinetic momenta satisfy
$[\pi_x,\pi_y]=\ii q\sh B_0$, giving
$E_\perp=\sh\omega_c(n_1+1/2)$,
$\omega_c=|qB_0|/m$.
For $qK>0$,
$H_z$ is an oscillator with
$\omega_z=\sqrt{2qK/m}$, so
\begin{equation}
\boxed{E_{n_1,n_2}=
\sh\omega_c(n_1+\tfrac12)+
\sh\omega_z(n_2+\tfrac12)}.
\end{equation}
The Landau levels have extensive transverse
degeneracy on an unbounded plane.
If $qK<0$, longitudinal motion is not
confined and the quoted discrete ladder does
not apply.

\problem{4.44}{Gauge covariance of the Schrödinger equation}
For real $\Lambda(\mathbf r,t)$, define
$\mathbf A'=\mathbf A+\nabla\Lambda$,
$\phi'=\phi-\partial_t\Lambda$,
and $\Psi'=\ee^{\ii q\Lambda/\sh}\Psi$.
Act on the new state with the
gauge-transformed covariant momentum:
\begin{equation}
(-\ii\sh\nabla-q\mathbf A')\Psi'
=\ee^{\ii q\Lambda/\sh}
(-\ii\sh\nabla-q\mathbf A)\Psi.
\end{equation}
Applying the operator again proves the same
identity for the kinetic-energy square.
The time derivative introduces
\begin{equation}
\ii\sh\partial_t\Psi'
=\ee^{\ii q\Lambda/\sh}
(\ii\sh\partial_t\Psi-
q\,\partial_t\Lambda\,\Psi).
\end{equation}
The extra term is precisely canceled by
$q\phi'=q(\phi-\partial_t\Lambda)$:
\begin{equation}
\boxed{\ii\sh\partial_t\Psi'=H'\Psi'}.
\end{equation}
Hence predictions are gauge invariant
provided the wavefunction also transforms.

\problem{4.45}{Energy shifts on a magnetic-flux-threaded ring}
A charge $q$ is constrained to a ring of
radius $b$ encircling magnetic flux $\Phi$.
Outside the solenoid $\mathbf B=0$, but
the circulating vector potential is
$A_\phi=\Phi/(2\pi b)$.
Minimal coupling gives
\begin{equation}
\boxed{H=\frac1{2mb^2}
\left(-\ii\sh\frac{\dd}{\dd\phi}
-\frac{q\Phi}{2\pi}\right)^2}.
\end{equation}
Periodic eigenfunctions
$\psi_n(\phi)=\ee^{\ii n\phi}/\sqrt{2\pi}$
yield
\begin{equation}
\boxed{E_n(\Phi)=\frac{\sh^2}{2mb^2}
\left(n-\frac{q\Phi}{2\pi\sh}\right)^2},
\qquad n\in\mathbb Z.
\end{equation}
The spectrum is periodic under
$\Phi\mapsto\Phi+\Phi_0$,
$\Phi_0=h/|q|$, after relabeling $n$.
On two distinct paths around the solenoid,
the relative phase is
$\boxed{\Delta\varphi=q\Phi/\sh}$,
which can shift an interference pattern
even when the particle never visits a region
with nonzero magnetic field.

\section{Additional three-dimensional methods (4.46--4.62)}
\problem{4.46}{Three-dimensional harmonic oscillator and degeneracies}
For $V(\mathbf r)=\frac12m\omega^2(x^2+y^2+z^2)$,
the Hamiltonian is the sum of three commuting
one-dimensional oscillator Hamiltonians.
Product eigenstates
$\psi_{n_x}(x)\psi_{n_y}(y)\psi_{n_z}(z)$
therefore have energies
\begin{equation}
E=\sh\omega(n_x+n_y+n_z+\tfrac32),
\qquad n_x,n_y,n_z=0,1,\dots.
\end{equation}
Writing $N=n_x+n_y+n_z$,
\begin{equation}
\boxed{E_N=(N+\tfrac32)\sh\omega}.
\end{equation}
The number of ordered nonnegative triples
whose entries sum to $N$ is
$\binom{N+2}{2}$ (stars and bars). Hence
\begin{equation}
\boxed{d(N)=\frac{(N+1)(N+2)}2}.
\end{equation}
The first degeneracies are $1,3,6,10,\dots$.

\problem{4.47}{Spherical-coordinate solution of the isotropic oscillator}
Separate $\psi=R(r)Y_\ell^m(\theta,\phi)$,
and introduce
$x=\sqrt{m\omega/\sh}\,r$ and
$\epsilon=2E/(\sh\omega)$.
The radial equation is
\begin{equation}
R''+\frac2xR'+
\left[\epsilon-x^2-\frac{\ell(\ell+1)}{x^2}\right]R=0.
\end{equation}
Correct near-origin and asymptotic behavior
suggests $R=x^\ell\ee^{-x^2/2}v(x)$.
After differentiation and cancellation,
\begin{equation}
v''+\left(\frac{2\ell+2}{x}-2x\right)v'
+[\epsilon-(2\ell+3)]v=0.
\end{equation}
Put $v(x)=\sum_{j=0}^\infty c_jx^j$.
The recurrence is
\begin{equation}
\boxed{c_{j+2}=
\frac{2j+2\ell+3-\epsilon}
{(j+2)(j+2\ell+3)}\,c_j}.
\end{equation}
Regularity selects an even series.
Normalizability requires truncation at
$j=2n_r$, which gives
$\epsilon=4n_r+2\ell+3$ and
\begin{equation}
\boxed{E_{n_r\ell}
=(2n_r+\ell+\tfrac32)\sh\omega},
\qquad n_r,\ell=0,1,\dots.
\end{equation}
Identify $N=2n_r+\ell$ from Problem 4.46.
For a given $N$, $\ell=N,N-2,\dots\ge0$.
Summing the $2\ell+1$ magnetic states yields
\begin{equation}
\boxed{\sum_{\ell=N,N-2,\dots}(2\ell+1)
=\frac{(N+1)(N+2)}2}.
\end{equation}
This checks complete agreement of
Cartesian and spherical separations.

\problem{4.48}{Three-dimensional virial theorem}
Take $G=\mathbf r\cdot\mathbf p$ and use
$\dd\avg G/\dd t=(\ii/\sh)\avg{[H,G]}$.
The canonical commutators give
\begin{equation}
\boxed{\frac{\dd\avg G}{\dd t}
=2\avg T-\avg{\mathbf r\cdot\nabla V}}.
\end{equation}
For a stationary state $\dd\avg G/\dd t=0$
and therefore
\begin{equation}
\boxed{2\avg T=\avg{\mathbf r\cdot\nabla V}}.
\end{equation}
For Coulomb $V\propto r^{-1}$,
$\mathbf r\cdot\nabla V=-V$,
hence $2\avg T=-\avg V$.
Using $E_n=\avg T+\avg V$ yields
\begin{equation}
\boxed{\avg T=-E_n,\qquad\avg V=2E_n}.
\end{equation}
For $V\propto r^2$,
$\mathbf r\cdot\nabla V=2V$, giving
$\boxed{\avg T=\avg V=E_n/2}$.
These identities assume normalizable
states and vanishing surface terms.

\problem{4.49}{Probability current and hydrogen orbital flow}
The three-dimensional current is
\begin{equation}
\boxed{\mathbf J=\frac{\sh}{2m\ii}
(\Psi^*\nabla\Psi-\Psi\nabla\Psi^*)}.
\end{equation}
Take the Schrödinger equation
$\ii\sh\partial_t\Psi=
-\sh^2\nabla^2\Psi/(2m)+V\Psi$,
multiply by $\Psi^*$, and subtract the
complex-conjugate equation multiplied by
$\Psi$. The real-potential terms cancel:
\begin{equation}
\partial_t|\Psi|^2=
-\nabla\cdot\frac{\sh}{2m\ii}
(\Psi^*\nabla\Psi-\Psi\nabla\Psi^*)
=\boxed{-\nabla\cdot\mathbf J}.
\end{equation}
Integrating over a fixed volume gives
$\oint\mathbf J\cdot\dd\mathbf a
=-\dd(\int_V|\Psi|^2\dd^3r)/\dd t$.
For $\psi_{211}=-r\sin\theta
\ee^{-r/(2\aB)}\ee^{\ii\phi}/
(8\sqrt\pi\aB^{5/2})$, only
the azimuthal phase contributes:
\begin{equation}
\mathbf J=\frac{\sh}{mr\sin\theta}
|\psi_{211}|^2\hat{\boldsymbol\phi}
=\boxed{\frac{\sh}{64\pi m\aB^5}
r\ee^{-r/\aB}\sin\theta\,
\hat{\boldsymbol\phi}}.
\end{equation}
The mass-flow orbital angular momentum is
$m\int(\mathbf r\times\mathbf J)\dd^3r$.
Its $z$ component is
$m\int r\sin\theta\,J_\phi\dd^3r=
\sh\int|\psi_{211}|^2\dd^3r$, so
\begin{equation}
\boxed{L_z=\sh},
\end{equation}
as expected for $m_\ell=+1$.

\problem{4.50}{Momentum-space hydrogen ground state}
For
$\psi_{100}(\mathbf r)=
\ee^{-r/\aB}/\sqrt{\pi\aB^3}$,
define the normalized Fourier transform
\begin{equation}
\phi(\mathbf p)=\frac{1}{(2\pi\sh)^{3/2}}
\int_{\mathbb R^3}\ee^{-\ii\mathbf p\cdot
\mathbf r/\sh}\psi_{100}(\mathbf r)\dd^3r.
\end{equation}
Choose the polar axis along $\mathbf p$.
The angular integral is
$\int\dd\Omega\,\ee^{-\ii pr\cos\theta/\sh}
=4\pi\sin(pr/\sh)/(pr/\sh)$.
Using
$\int_0^\infty r\ee^{-\beta r}\sin(kr)\dd r
=2\beta k/(\beta^2+k^2)^2$, gives
\begin{equation}
\boxed{\phi(\mathbf p)=\frac1\pi
\left(\frac{2\aB}{\sh}\right)^{3/2}
\frac1{[1+(\aB p/\sh)^2]^2}}.
\end{equation}
Put $u=\aB p/\sh$. Its norm is
\begin{equation}
\int|\phi|^2\dd^3p
=\frac{32}{\pi}\int_0^\infty
\frac{u^2\dd u}{(1+u^2)^4}=1,
\end{equation}
where the integral equals $\pi/32$.
Similarly,
\begin{equation}
\avg{p^2}=\frac{32\sh^2}{\pi\aB^2}
\int_0^\infty\frac{u^4\dd u}{(1+u^2)^4}
=\boxed{\frac{\sh^2}{\aB^2}}.
\end{equation}
Thus $\boxed{\avg T=\sh^2/(2m\aB^2)=-E_1}$,
in agreement with the Coulomb virial theorem.

\problem{4.51}{Why every attractive circular well binds in two dimensions}
For a radially symmetric two-dimensional
potential $V=-V_0$ inside $r<a$
and zero outside, a bound $s$-wave has
$-V_0<E<0$, with
$k^2=2m(V_0+E)/\sh^2$ and
$\kappa^2=-2mE/\sh^2$.
The finite-at-origin interior solution is
$R_{\rm in}=AJ_0(kr)$; the decaying exterior
solution is $R_{\rm out}=BK_0(\kappa r)$.
Matching $R$ and $R'$ at $r=a$ gives
\begin{equation}
\boxed{k\frac{J_1(ka)}{J_0(ka)}
=\kappa\frac{K_1(\kappa a)}{K_0(\kappa a)}}.
\end{equation}
For arbitrarily weak attraction, $ka\ll1$.
As $\kappa a\to0$,
$K_0(\kappa a)\sim-\ln(\kappa a/2)-\gamma_E$,
$\kappa K_1(\kappa a)\sim1/a$, while
$kJ_1(ka)/J_0(ka)\sim k^2a/2$.
Hence
\begin{equation}
-\ln(\kappa a/2)-\gamma_E
\sim\frac{2}{(ka)^2},
\quad
\boxed{\kappa a\sim
2\exp[-\gamma_E-2/(ka)^2]}.
\end{equation}
For every $V_0>0$ this is nonzero,
though exponentially small, so there
is at least one negative-energy bound state.
Unlike three dimensions, there is no
nonzero threshold depth for an attractive
two-dimensional circular well.

\problem{4.52}{Hydrogen $\psi_{321}$ and radial moments}
For $n=3,\ell=2$,
$R_{32}=(2/(3\aB))^{3/2}\rho^2
\ee^{-\rho/2}/\sqrt{720}$,
$\rho=2r/(3\aB)$.
With
$Y_2^1=-\sqrt{15/(8\pi)}
\sin\theta\cos\theta\ee^{\ii\phi}$,
simplification gives
\begin{equation}
\boxed{\psi_{321}(r,\theta,\phi)=
-\frac{r^2}{81\sqrt\pi\,\aB^{7/2}}
\ee^{-r/(3\aB)}
\sin\theta\cos\theta\ee^{\ii\phi}}.
\end{equation}
For normalization, the angular integral is
$\int\sin^2\theta\cos^2\theta\dd\Omega=8\pi/15$;
the radial integral is
\begin{equation}
\int_0^\infty r^6\ee^{-2r/(3\aB)}\dd r
=6!\left(\frac{3\aB}{2}\right)^7.
\end{equation}
Combining with the displayed prefactor yields one.
For arbitrary real $s$,
\begin{equation}
\boxed{\avg{r^s}=
\frac{\Gamma(s+7)}{\Gamma(7)}
\left(\frac{3\aB}{2}\right)^s},
\qquad\boxed{s>-7}.
\end{equation}
The restriction follows from the
$r^{s+6}$ behavior near $r=0$.

\problem{4.53}{Hydrogen $\psi_{433}$ and transverse angular momentum}
For $n=4,\ell=3,m=3$,
$R_{43}=(1/(2\aB))^{3/2}\rho^3
\ee^{-\rho/2}/\sqrt{8\,7!}$,
where $\rho=r/(2\aB)$, and
$Y_3^3=-\sqrt{35/(64\pi)}
\sin^3\theta\ee^{3\ii\phi}$.
Therefore
\begin{equation}
\boxed{\psi_{433}=
-\frac{r^3}{6144\sqrt\pi\,\aB^{9/2}}
\ee^{-r/(4\aB)}
\sin^3\theta\ee^{3\ii\phi}}.
\end{equation}
The exact hydrogen radial expectation is
$\avg r=\frac{\aB}{2}[3n^2-\ell(\ell+1)]$;
here
\begin{equation}
\boxed{\avg r=18\aB}.
\end{equation}
Since $L_x^2+L_y^2=L^2-L_z^2$,
the state is already an eigenstate of the
requested observable:
\begin{equation}
\boxed{(L_x^2+L_y^2)\psi_{433}
=\sh^2[3(4)-3^2]\psi_{433}
=3\sh^2\psi_{433}}.
\end{equation}
A measurement returns $3\sh^2$ with probability one.

\problem{4.54}{Ground-state penetration of a finite-size nucleus}
For the hydrogen $1s$ state,
the probability inside a sphere of radius $b$ is
\begin{align}
P(r<b)&=\frac4{\aB^3}\int_0^b
r^2\ee^{-2r/\aB}\dd r\\
&=\boxed{1-\ee^{-2b/\aB}
\left(1+\frac{2b}{\aB}+
\frac{2b^2}{\aB^2}\right)}.
\end{align}
Let $\epsilon=2b/\aB\ll1$.
Expanding
$1-\ee^{-\epsilon}(1+\epsilon+\epsilon^2/2)
=\epsilon^3/6+O(\epsilon^4)$ gives
\begin{equation}
\boxed{P\simeq\frac43(b/\aB)^3}.
\end{equation}
Approximating $\psi$ by its value at the
origin in the nuclear volume gives
$(4\pi b^3/3)|\psi(0)|^2=
(4/3)(b/\aB)^3$, the same result.
With $b=10^{-15}$ m and
$\aB=0.5\times10^{-10}$ m,
\begin{equation}
\boxed{P\simeq1.07\times10^{-14}}.
\end{equation}
This uses a point-Coulomb wavefunction
inside the nuclear volume, as requested;
finite nuclear structure modifies it.

\problem{4.55}{Circular hydrogenic states and radial uncertainty}
For $\ell=n-1$, the radial polynomial
is constant, and
$R_{n,n-1}=N_n r^{n-1}\ee^{-r/(n\aB)}$.
Set $\beta=2/(n\aB)$.
Normalization requires
$1=|N_n|^2\int_0^\infty r^{2n}
\ee^{-\beta r}\dd r=
|N_n|^2(2n)!/\beta^{2n+1}$, so
\begin{equation}
\boxed{N_n=\frac{\beta^{n+1/2}}{\sqrt{(2n)!}}}.
\end{equation}
The radial moments follow from successive
gamma-function integrals:
\begin{align}
\boxed{\avg r}&=\frac{2n+1}{\beta}
=\frac{n\aB}{2}(2n+1),\\
\boxed{\avg{r^2}}&=\frac{(2n+1)(2n+2)}{\beta^2}.
\end{align}
Thus
$\sigma_r^2=(2n+1)/\beta^2$,
and
\begin{equation}
\boxed{\sigma_r=\frac{\avg r}{\sqrt{2n+1}}}.
\end{equation}
The fractional radial spread decreases
as $n^{-1/2}$ for large $n$,
suggesting a more localized radial
distribution in highly excited circular states.

\problem{4.56}{Two distinct hydrogen transitions with the same wavelength}
The Rydberg formula assigns
$\lambda^{-1}=R_H(n_f^{-2}-n_i^{-2})$.
Take the two different transitions
$(n_i,n_f)=(7,5)$ and $(35,7)$:
\begin{align}
\frac1{5^2}-\frac1{7^2}
&=\frac{49-25}{1225}=\frac{24}{1225},\\
\frac1{7^2}-\frac1{35^2}
&=\frac{25-1}{1225}=\frac{24}{1225}.
\end{align}
Both have $n_i>n_f$, hence
\begin{equation}
\boxed{\lambda_{7\to5}=\lambda_{35\to7}
=\frac{1225}{24R_H}}.
\end{equation}
This pair is not the example given
in the textbook.

\problem{4.57}{Uncertainty between $x^2$ and $L_z$}
Using $[L_z,x]=\ii\sh y$, one obtains
$[x^2,L_z]=-\ii\sh(xy+yx)
=-2\ii\sh xy$,
because $x$ and $y$ commute.
Robertson's inequality gives
\begin{equation}
\boxed{\sigma_{x^2}\,\sigma_{L_z}
\ge\sh|\avg{xy}|}.
\end{equation}
Every $\psi_{n\ell m}$ is an $L_z$
eigenstate, so $\sigma_{L_z}=0$.
Consequently
\begin{equation}
\boxed{\avg{xy}_{n\ell m}=0}.
\end{equation}
This agrees with azimuthal symmetry
of $|\psi_{n\ell m}|^2$.

\problem{4.58}{Probabilities for a complex spinor}
Take $\chi=A(1-2\ii,2)^T$.
Normalization gives $|A|^2(5+4)=1$;
choose $\boxed{A=1/3}$.
For $S_z$, the outcomes are
$\pm\sh/2$ with
\begin{equation}
\boxed{P_z(+)={5}/{9},\quad
P_z(-)={4}/{9},\quad
\avg{S_z}=\sh/18}.
\end{equation}
Since $a^*b=(1+2\ii)2/9=(2+4\ii)/9$,
$\avg{S_x}=\sh\,\mathrm{Re}(a^*b)=2\sh/9$,
$\avg{S_y}=\sh\,\mathrm{Im}(a^*b)=4\sh/9$.
The two-outcome relations
$P_i(\pm)=\frac12\pm\avg{S_i}/\sh$ give
\begin{align}
\boxed{P_x(+)=13/18,\ P_x(-)=5/18},\\
\boxed{P_y(+)=17/18,\ P_y(-)=1/18}.
\end{align}
Each pair sums to one and the
expectations follow from their weighted outcomes.

\problem{4.59}{Opposite spin correlations in the singlet}
The singlet is
$\ket{\mathrm s}=
(\ket{\uparrow\downarrow}
-\ket{\downarrow\uparrow})/\sqrt2$.
It satisfies
$\bra{\mathrm s}\sigma_i^{(1)}
\sigma_j^{(2)}\ket{\mathrm s}=-\delta_{ij}$:
the diagonal terms are $-1$ by rotational
invariance, and mixed terms vanish.
For unit vectors $\mathbf a,\mathbf b$,
\begin{align}
\avg{S_{\mathbf a}^{(1)}
S_{\mathbf b}^{(2)}}&=
\frac{\sh^2}{4}
\sum_{ij}a_ib_j
\avg{\sigma_i^{(1)}\sigma_j^{(2)}}\\
&=-\frac{\sh^2}{4}\,\mathbf a\cdot\mathbf b.
\end{align}
Writing $\mathbf a\cdot\mathbf b=\cos\theta$ gives
\begin{equation}
\boxed{\avg{S_{\mathbf a}^{(1)}
S_{\mathbf b}^{(2)}}=-\frac{\sh^2}{4}\cos\theta}.
\end{equation}

\problem{4.60}{Coupling spin one-half to arbitrary spin}
Let $s_2=j$ and set total $s=j\pm1/2$.
At fixed $m$, only two uncoupled states occur:
$\ket{+\frac12,j,m-\frac12}$ and
$\ket{-\frac12,j,m+\frac12}$.
Write
$\ket{sm}=A\ket{+\frac12,j,m-\frac12}
+B\ket{-\frac12,j,m+\frac12}$.
The matrix for $S^2$ on this span is
obtained from
$S^2=S_1^2+S_2^2+2S_{1z}S_{2z}+
S_{1+}S_{2-}+S_{1-}S_{2+}$.
Its diagonalization, followed by
$|A|^2+|B|^2=1$, yields
\begin{align}
\boxed{s=j+\tfrac12:}&\quad
A=\sqrt{\frac{j+m+1/2}{2j+1}},&
B=\sqrt{\frac{j-m+1/2}{2j+1}},\\
\boxed{s=j-\tfrac12:}&\quad
A=\sqrt{\frac{j-m+1/2}{2j+1}},&
B=-\sqrt{\frac{j+m+1/2}{2j+1}}.
\end{align}
This is the standard Condon--Shortley phase choice.
For $j=1/2$, these reproduce the triplet
and singlet coefficients; for $j=1$ and
$j=3/2$, substitution reproduces
the corresponding table entries.

\problem{4.61}{Spin-three-halves $S_x$ matrix and eigenvalues}
In the ordered $S_z$ basis
$m=3/2,1/2,-1/2,-3/2$,
\begin{equation}
\boxed{S_x=\frac\sh2
\begin{pmatrix}
0&\sqrt3&0&0\\
\sqrt3&0&2&0\\
0&2&0&\sqrt3\\
0&0&\sqrt3&0
\end{pmatrix}}.
\end{equation}
Writing $t=2\lambda/\sh$, the
characteristic polynomial is
\begin{equation}
\det(tI-2S_x/\sh)
=(t^2-1)(t^2-9).
\end{equation}
Therefore
\begin{equation}
\boxed{\lambda=\pm\frac{3\sh}{2},
\quad\pm\frac\sh2}.
\end{equation}
These eigenvalues agree with rotational
equivalence of $S_x$ and $S_z$.

\problem{4.62}{General spin-$s$ Cartesian matrices}
In the ordered $\ket{s,m}$ basis
$m=s,s-1,\ldots,-s$,
\begin{align}
\boxed{(S_z)_{m'm}}&=
\sh m\,\delta_{m'm},\\
\boxed{(S_+)_{m'm}}&=
\sh\sqrt{(s-m)(s+m+1)}\,
\delta_{m',m+1},\\
\boxed{(S_-)_{m'm}}&=
\sh\sqrt{(s+m)(s-m+1)}\,
\delta_{m',m-1}.
\end{align}
Hence
\begin{equation}
\boxed{S_x=\tfrac12(S_++S_-),\qquad
S_y=\tfrac1{2\ii}(S_+-S_-)}.
\end{equation}
They are $(2s+1)\times(2s+1)$ Hermitian
matrices with nonzero off-diagonal entries
only between adjacent $m$ values.
Their commutators obey
$[S_i,S_j]=\ii\sh\epsilon_{ijk}S_k$
and $S^2=\sh^2s(s+1)I$.

\section{Advanced angular momentum and gauge problems (4.63--4.76)}
\problem{4.63}{Normalization of all spherical harmonics}
Write $Y_\ell^m=K_{\ell m}
P_\ell^m(\cos\theta)\ee^{\ii m\phi}$.
The differential ladder operator and
the identity for associated Legendre
functions imply
\begin{equation}
L_+Y_\ell^m=\sh K_{\ell m}
P_\ell^{m+1}(\cos\theta)
\ee^{\ii(m+1)\phi}.
\end{equation}
But the normalized angular-momentum rule is
$L_+Y_\ell^m=
\sh\sqrt{(\ell-m)(\ell+m+1)}Y_\ell^{m+1}$.
Consequently
\begin{equation}
\boxed{\frac{K_{\ell,m+1}}{K_{\ell m}}
=\frac1{\sqrt{(\ell-m)(\ell+m+1)}}}.
\end{equation}
Iterating from $m=0$ gives
\begin{equation}
K_{\ell m}=K_{\ell0}
\sqrt{\frac{(\ell-m)!}{(\ell+m)!}}.
\end{equation}
Orthogonality of ordinary Legendre polynomials
requires
$1=2\pi K_{\ell0}^2\int_{-1}^1P_\ell^2\dd x
=4\pi K_{\ell0}^2/(2\ell+1)$.
Thus
\begin{equation}
\boxed{K_{\ell m}=
\sqrt{\frac{2\ell+1}{4\pi}
\frac{(\ell-m)!}{(\ell+m)!}}}.
\end{equation}
The phase is fixed by the Condon--Shortley
convention built into $P_\ell^m$.
Negative $m$ follows from
$Y_\ell^{-m}=(-1)^mY_\ell^{m*}$.

\problem{4.64}{Combined orbital and spin hydrogen measurements}
Consider the normalized state
\begin{equation}
\ket\Psi=\sqrt{\frac13}\ket{\ell=1,m_\ell=0}
\ket{\uparrow}+
\sqrt{\frac23}\ket{1,1}\ket{\downarrow},
\end{equation}
with the common normalized $R_{21}(r)$.
Both components have $\ell=1$ and $s=1/2$.
Therefore
\begin{equation}
\boxed{L^2=2\sh^2\;(P=1),\qquad
S^2=\tfrac34\sh^2\;(P=1)}.
\end{equation}
For $L_z$ the outcomes are $0$ with $P=1/3$
and $\sh$ with $P=2/3$; for $S_z$, the outcomes
$\sh/2,-\sh/2$ have $P=1/3,2/3$.
Both components have $m_j=m_\ell+m_s=1/2$,
so $\boxed{J_z=\sh/2\;(P=1)}$.
For $j=3/2,m_j=1/2$ the CG vector is
$\sqrt{2/3}\ket{0,\uparrow}+
\sqrt{1/3}\ket{1,\downarrow}$.
Projecting gives amplitude
$\sqrt2/3+\sqrt2/3=2\sqrt2/3$, hence
\begin{equation}
\boxed{P(J^2=\tfrac{15}{4}\sh^2)=\tfrac89,\qquad
P(J^2=\tfrac34\sh^2)=\tfrac19}.
\end{equation}
For position measurement the spin alternatives
are orthogonal, so no spatial interference term:
\begin{align}
|\Psi(\mathbf r)|^2&=|R_{21}(r)|^2
\left[\tfrac13|Y_1^0|^2+
\tfrac23|Y_1^1|^2\right]\\
&=\boxed{\frac{|R_{21}(r)|^2}{4\pi}},
\end{align}
which is isotropic. The probability per unit
radius for spin up is
\begin{equation}
\boxed{P_+(r)=\frac13r^2|R_{21}(r)|^2}.
\end{equation}

\problem{4.65}{Three spin halves: one quadruplet and two doublets}
Use the shorthand $u=\uparrow$, $d=\downarrow$.
The totally symmetric spin-$3/2$ quadruplet is
\begin{align}
\ket{\tfrac32,\tfrac32}&=\ket{uuu},\\
\ket{\tfrac32,\tfrac12}&=
(\ket{uud}+\ket{udu}+\ket{duu})/\sqrt3,\\
\ket{\tfrac32,-\tfrac12}&=
(\ket{udd}+\ket{dud}+\ket{ddu})/\sqrt3,\\
\ket{\tfrac32,-\tfrac32}&=\ket{ddd}.
\end{align}
The first spin-$1/2$ doublet is obtained by
placing particles 1 and 2 in a singlet:
\begin{align}
\ket{\tfrac12,\tfrac12}_{1}
&=(\ket{udu}-\ket{duu})/\sqrt2,\\
\ket{\tfrac12,-\tfrac12}_{1}
&=(\ket{udd}-\ket{dud})/\sqrt2.
\end{align}
A second independent doublet, orthogonal
to both the first doublet and the quadruplet,
is
\begin{align}
\ket{\tfrac12,\tfrac12}_{2}
&=(2\ket{uud}-\ket{udu}-\ket{duu})/\sqrt6,\\
\ket{\tfrac12,-\tfrac12}_{2}
&=(\ket{udd}+\ket{dud}-2\ket{ddu})/\sqrt6.
\end{align}
Each state has unit norm, the two doublets
are mutually orthogonal, and total
ladder operators connect the $\pm1/2$
members within each doublet. Dimension
count gives $4+2+2=8=2^3$.

\problem{4.66}{Saturation of the spin uncertainty bound}
Let $\chi=(a,b)^T$, $\|\chi\|=1$.
Remove an overall phase to make $a\ge0$ real.
Define the Bloch components
$n_x=2a\operatorname{Re}b$,
$n_y=2a\operatorname{Im}b$,
$n_z=|a|^2-|b|^2$.
A pure spinor obeys $n_x^2+n_y^2+n_z^2=1$.
Since $\sigma_{S_i}^2=
(\sh^2/4)(1-n_i^2)$,
\begin{align}
\sigma_{S_x}^2\sigma_{S_y}^2
-\frac{\sh^2}{4}\avg{S_z}^2
&=\frac{\sh^4}{16}
[(1-n_x^2)(1-n_y^2)-n_z^2]\\
&=\frac{\sh^4}{16}n_x^2n_y^2.
\end{align}
Equality therefore holds if and only if
$n_xn_y=0$. With $a$ real, this means
\begin{equation}
\boxed{\operatorname{Re}b=0
\quad\text{or}\quad\operatorname{Im}b=0},
\end{equation}
so $b$ can be chosen purely imaginary
or purely real. The special case $a=0$
automatically saturates the inequality.

\problem{4.67}{Antiferromagnetic frustration on triangles and squares}
For three spins $\mathbf S_i$ of magnitude
$s_i=1/2$, define $\mathbf S=\sum_i\mathbf S_i$.
Then
\begin{equation}
S^2=\sum_iS_i^2+
2(S_1\cdot S_2+S_2\cdot S_3+S_3\cdot S_1).
\end{equation}
Hence for the triangular Hamiltonian
\begin{equation}
\boxed{H_\triangle=
\frac J2(S^2-\tfrac94\sh^2)}.
\end{equation}
Possible total spins are $3/2$ and $1/2$;
the latter occurs in two independent doublets.
Their energies are
$E_{1/2}=-3J\sh^2/4$ and
$E_{3/2}=3J\sh^2/4$.
Thus
\begin{equation}
\boxed{E_{\rm ground}^{\triangle}=-\tfrac34J\sh^2,
\qquad d_{\rm ground}=4}.
\end{equation}
For a square, group opposite vertices:
$\mathbf A=\mathbf S_1+\mathbf S_3$,
$\mathbf B=\mathbf S_2+\mathbf S_4$.
Then $H_\square=J\mathbf A\cdot\mathbf B$
and
\begin{equation}
\boxed{H_\square=
\frac J2(S^2-A^2-B^2)}.
\end{equation}
Each opposite-vertex pair can have
spin $0$ or $1$. Choosing $A=B=1$
and combining to total $S=0$ minimizes
the energy:
\begin{equation}
\boxed{E_{\rm ground}^{\square}=-2J\sh^2},
\end{equation}
with a unique singlet ground state.

\problem{4.68}{Numerical radial hydrogen in a large spherical wall}
For a Coulomb center with a hard wall
$u(b)=0$, discretize
$u(r)=rR(r)$ at $r_j=j\Delta r$,
$\Delta r=b/(N+1)$, and put
$\beta=\Delta r/\aB=b/[(N+1)\aB]$.
With energies scaled to $\sh^2/(2m\Delta r^2)$,
the interior tridiagonal matrix has
\begin{equation}
\boxed{H_{jj}=2+v_j,\quad
H_{j,j\pm1}=-1,\quad
v_j=-\frac{2\beta}{j}
+\frac{\ell(\ell+1)}{j^2}}.
\end{equation}
The end conditions are $u_0=u_{N+1}=0$.
For $\beta=1/50$, $N=1000$,
$b=(N+1)\aB/50\simeq20.02\aB$.
Diagonalizing the symmetric tridiagonal matrix,
and converting to Hartrees
$(1\ {\rm Ha}=\sh^2/(m\aB^2))$, gives
\begin{equation}
\begin{array}{c|rrr}
\ell& E_1/{\rm Ha}&E_2/{\rm Ha}&E_3/{\rm Ha}\\\hline
0&-0.499950&-0.124984&-0.049955\\
1&-0.124996&-0.051639&+0.007931\\
2&-0.053979&-0.005638&+0.067981
\end{array}.
\end{equation}
The corresponding unconfined hydrogen energies
are $E_n=-1/(2n^2)$ Ha, with
$n\ge\ell+1$. Higher-energy states are
strongly affected by the finite wall; in
particular positive confined eigenvalues
are not hydrogen bound levels.
The attached numerical script reproduces
the matrix eigenvalues and radial plots.
\note{The geometry here is a spherical \emph{volume}
with a hard radial boundary, not a surface-constrained
particle or a toroidal hydrogen system.}

\problem{4.69}{Numerical shooting targets for the first Coulomb states}
For the hydrogen radial equation, set
$E=-\sh^2/(2m\aB^2 n_{\rm eff}^2)$
and $\rho=2r/(n_{\rm eff}\aB)$.
The radial ODE is
\begin{equation}
\frac{\dd^2u}{\dd\rho^2}
+\left[-\frac14+\frac{n_{\rm eff}}{\rho}
-\frac{\ell(\ell+1)}{\rho^2}\right]u=0.
\end{equation}
Near $\rho=0$ the regular solution has
$u\propto\rho^{\ell+1}$.
Integrate from a small positive cutoff
$\rho_{\min}$ with
$u(\rho_{\min})=\rho_{\min}^{\ell+1}$
and $u'(\rho_{\min})=
(\ell+1)\rho_{\min}^\ell$.
For varying $n_{\rm eff}$, demand a
decaying rather than growing tail as
$\rho\to\infty$. The first three regular
normalizable values are
\begin{equation}
\boxed{\begin{array}{c|ccc}
\ell&n_{\rm eff,1}&n_{\rm eff,2}&n_{\rm eff,3}\\\hline
0&1.000&2.000&3.000\\
1&2.000&3.000&4.000\\
2&3.000&4.000&5.000
\end{array}}
\end{equation}
to four significant figures.
These values follow analytically from
termination of the confluent-hypergeometric
series at $n\ge\ell+1$. Shooting is
a numerical independent check; convergence
must be tested with larger integration
cutoffs and smaller step sizes.

\problem{4.70}{Sequential Stern--Gerlach measurements}
(a) For an initial $\ket{z+}$,
a first $z$ measurement always yields $z+$.
If the second measurement retains only
$x+$, every retained particle had $z+$
earlier. If instead the first measurement
was along $x$, the $x+$ and $x-$ outcomes
each occur with probability $1/2$,
but the second $x$ measurement retains only
particles that were previously $x+$.
Thus in either version
$\boxed{P(\text{first up}\mid\text{second }x+)=1}$.
(b) For arbitrary unit axes $\mathbf a,\mathbf b,
\mathbf c$, let
$u=\cos^2(\theta_{ab}/2)$,
$v=\cos^2(\theta_{bc}/2)$.
The joint probabilities for the two successive
measurements along $\mathbf b$ and $\mathbf c$
are
\begin{align}
P_{++}&=uv,&P_{+-}&=u(1-v),\\
P_{-+}&=(1-u)(1-v),&
P_{--}&=(1-u)v.
\end{align}
Among the subensemble whose second result is $c+$,
the fraction with an earlier $b+$ is
\begin{equation}
\boxed{\frac{uv}{uv+(1-u)(1-v)}
=\frac1{1+\tan^2(\theta_{ab}/2)
\tan^2(\theta_{bc}/2)}}.
\end{equation}

\problem{4.71}{Cartesian hydrogen $2p$ orbitals}
The combinations of spherical harmonics satisfy
\begin{align}
\sin\theta\cos\phi\,\sqrt{\frac{3}{4\pi}}
&=(Y_1^{-1}-Y_1^1)/\sqrt2,\\
\sin\theta\sin\phi\,\sqrt{\frac{3}{4\pi}}
&=\ii(Y_1^{-1}+Y_1^1)/\sqrt2,\\
\cos\theta\,\sqrt{\frac{3}{4\pi}}&=Y_1^0.
\end{align}
Multiplying by the common normalized $R_{21}(r)$,
\begin{equation}
\boxed{\psi_{2p_x}=
(\psi_{21,-1}-\psi_{211})/\sqrt2,\quad
\psi_{2p_y}=
\ii(\psi_{21,-1}+\psi_{211})/\sqrt2,\quad
\psi_{2p_z}=\psi_{210}}.
\end{equation}
These are proportional to
$x\ee^{-r/(2\aB)}$, $y\ee^{-r/(2\aB)}$,
$z\ee^{-r/(2\aB)}$, respectively.
Since, for example,
$L_x=-\ii\sh(y\partial_z-z\partial_y)$
and it annihilates any function
$x f(r)$, one has
\begin{equation}
\boxed{L_x\psi_{2p_x}=0,\quad
L_y\psi_{2p_y}=0,\quad
L_z\psi_{2p_z}=0}.
\end{equation}
Each aligned orbital is an eigenstate of
the corresponding angular-momentum component
with eigenvalue zero. Isodensity contours have
two lobes along the appropriate Cartesian axis,
separated by its perpendicular nodal plane.

\problem{4.72}{Uniform-field vector potential and magnetic coupling}
Let $\mathbf A=-\frac12\mathbf r\times\mathbf B_0
=\frac12\mathbf B_0\times\mathbf r$,
with constant $\mathbf B_0$.
The vector identity
$\nabla\times(\mathbf B_0\times\mathbf r)=
2\mathbf B_0$ gives
\begin{equation}
\boxed{\nabla\times\mathbf A=\mathbf B_0}.
\end{equation}
For a charge $q$ with spin coupling
$-\gamma\mathbf B_0\cdot\mathbf S$,
\begin{equation}
H=\frac{(\mathbf p-q\mathbf A)^2}{2m}
+q\phi-\gamma\mathbf B_0\cdot\mathbf S.
\end{equation}
Since $\nabla\cdot\mathbf A=0$,
$\mathbf p\cdot\mathbf A=\mathbf A\cdot\mathbf p$.
Also
$\mathbf A\cdot\mathbf p=
\frac12\mathbf B_0\cdot(\mathbf r\times\mathbf p)
=\frac12\mathbf B_0\cdot\mathbf L$
and
$\mathbf A^2=\frac14
[r^2B_0^2-(\mathbf r\cdot\mathbf B_0)^2]$.
Thus
\begin{equation}
\boxed{
H=\frac{p^2}{2m}+q\phi
-\mathbf B_0\cdot\left(\frac q{2m}\mathbf L+
\gamma\mathbf S\right)
+\frac{q^2}{8m}
[r^2B_0^2-(\mathbf r\cdot\mathbf B_0)^2]}.
\end{equation}
The term linear in field gives
the orbital/spin Zeeman interaction;
the quadratic term describes
diamagnetic coupling.

\problem{4.73}{Quantum force in a Stern--Gerlach field}
For a neutral spin-one-half particle
$H=p^2/(2m)-\gamma\mathbf B(\mathbf r)\cdot\mathbf S$.
With
$\mathbf B=-\alpha x\hat x+
(B_0+\alpha z)\hat z$,
only the spin interaction fails
to commute with $p_z$:
\begin{align}
\frac{\dd\avg{p_z}}{\dd t}
&=\frac{\ii}{\sh}
\avg{[H,p_z]}
=\gamma\avg{\partial_z\mathbf B\cdot\mathbf S}\\
&=\gamma\alpha\avg{S_z}.
\end{align}
Since $\dd\avg z/\dd t=\avg{p_z}/m$,
\begin{equation}
\boxed{m\frac{\dd^2\avg z}{\dd t^2}
=\gamma\alpha\avg{S_z}}.
\end{equation}
This is the quantum expectation-value
version of the Stern--Gerlach force law.

\problem{4.74}{Coupled spinor Schrödinger equations for Stern--Gerlach}
Write the complete state as
$\Psi=\Psi_+\chi_++\Psi_-\chi_-$,
and use
$\mathbf B=(-\alpha x,0,B_0+\alpha z)$.
The Hamiltonian matrix is
\begin{equation}
H=-\frac{\sh^2}{2m}\nabla^2I
-\frac{\gamma\sh}{2}
\begin{pmatrix}
B_0+\alpha z&-\alpha x\\
-\alpha x&-B_0-\alpha z
\end{pmatrix}.
\end{equation}
Equating spin components gives
\begin{align}
\ii\sh\partial_t\Psi_+
&=-\frac{\sh^2}{2m}\nabla^2\Psi_+
-\frac{\gamma\sh}{2}(B_0+\alpha z)\Psi_+
+\frac{\gamma\sh\alpha x}{2}\Psi_-,\\
\ii\sh\partial_t\Psi_-
&=-\frac{\sh^2}{2m}\nabla^2\Psi_-
+\frac{\gamma\sh}{2}(B_0+\alpha z)\Psi_-
+\frac{\gamma\sh\alpha x}{2}\Psi_+.
\end{align}
Factor the rapid uniform-field precession:
$\Psi_\pm=\ee^{\pm\ii\gamma B_0t/2}
\widetilde\Psi_\pm$. The equations become
\begin{align}
\ii\sh\partial_t\widetilde\Psi_+
&=\left[-\frac{\sh^2}{2m}\nabla^2
-\frac{\gamma\sh\alpha z}{2}\right]\widetilde\Psi_+
+\frac{\gamma\sh\alpha x}{2}
\ee^{-\ii\gamma B_0t}\widetilde\Psi_-,\\
\ii\sh\partial_t\widetilde\Psi_-
&=\left[-\frac{\sh^2}{2m}\nabla^2
+\frac{\gamma\sh\alpha z}{2}\right]\widetilde\Psi_-
+\frac{\gamma\sh\alpha x}{2}
\ee^{+\ii\gamma B_0t}\widetilde\Psi_+.
\end{align}
When rapid off-diagonal oscillations are
negligible, the spin components independently
see the linear potentials
$V_\pm=\mp\gamma\sh\alpha z/2$.
Their mean accelerations are
\begin{equation}
\boxed{\ddot{\avg z}_\pm
=\pm\frac{\gamma\sh\alpha}{2m}}.
\end{equation}
The packet splits into two paths.
\note{Neglecting off-diagonal coupling is
a controlled approximation only when the
uniform-field precession is rapid compared
with the spatial coupling dynamics.}

\problem{4.75}{Time-dependent magnetic flux in the ring problem}
For charge $q$ on a ring of radius $b$,
the time-dependent Hamiltonian is
\begin{equation}
H(t)=\frac{\sh^2}{2mb^2}
\left(-\ii\frac{\dd}{\dd\phi}
-\frac{q\Phi(t)}{2\pi\sh}\right)^2.
\end{equation}
All $H(t)$ are functions of the
same derivative operator and commute
at different times. A Fourier angular
mode $\ee^{\ii n\phi}/\sqrt{2\pi}$
has instantaneous eigenvalue
\begin{equation}
E_n(t)=\frac{\sh^2}{2mb^2}
\left(n-\frac{q\Phi(t)}{2\pi\sh}\right)^2.
\end{equation}
Let $f(t)=\sh^{-1}\int_0^tE_n(t')\dd t'$.
Then
\begin{equation}
\boxed{\Psi_n(\phi,t)=
\frac{\ee^{\ii n\phi}}{\sqrt{2\pi}}
\ee^{-\ii f(t)}}
\end{equation}
obeys
$\ii\sh\partial_t\Psi_n=
\sh f'(t)\Psi_n=E_n(t)\Psi_n=H(t)\Psi_n$.
Normalization is unchanged, as only
a time-dependent global phase is acquired.

\problem{4.76}{Mechanical angular momentum and changing flux}
Faraday's law gives the induced azimuthal
electric field on the ring:
$2\pi bE_\phi=-\dd\Phi/\dd t$.
For tangential speed $v=b\omega$,
the electric work rate is
\begin{equation}
\frac{\dd W}{\dd t}=qvE_\phi
=-\frac{q\omega}{2\pi}\frac{\dd\Phi}{\dd t},
\quad\boxed{\frac{\dd W}{\dd\Phi}
=-\frac{q\omega}{2\pi}}.
\end{equation}
In a state $\psi_n=\ee^{\ii n\phi}/\sqrt{2\pi}$,
canonical angular momentum is
$L_{z,\mathrm{can}}=n\sh$.
With $A_\phi=\Phi/(2\pi b)$, the
mechanical momentum is
$p_{\phi,\mathrm{mech}}=n\sh/b-qA_\phi$,
so
\begin{equation}
\boxed{L_{z,\mathrm{mech}}=
n\sh-\frac{q\Phi}{2\pi}},
\quad
\boxed{\omega=\frac{n\sh-q\Phi/(2\pi)}{mb^2}}.
\end{equation}
Differentiating the quantum stationary
energy at fixed integer $n$,
\begin{align}
\frac{\dd E_n}{\dd\Phi}
&=\frac{\dd}{\dd\Phi}
\frac{[n\sh-q\Phi/(2\pi)]^2}{2mb^2}\\
&=-\frac{q}{2\pi mb^2}
[n\sh-q\Phi/(2\pi)]\\
&=\boxed{-\frac{q\omega}{2\pi}}.
\end{align}
This exactly agrees with the
classical work rate per unit flux.
The \emph{mechanical} angular momentum
is not generally an integer multiple
of $\sh$ even though canonical
angular momentum is.


\section{Reproducible numerical illustrations}
The accompanying \texttt{scripts/verify.py} checks a set of independent
eigenvalue, normalization, matrix, and spin-coupling calculations.
\IfFileExists{figures/radial_density.pdf}{
\begin{figure}[htbp]\centering
\includegraphics[width=.86\linewidth]{figures/radial_density.pdf}
\caption{Hydrogenic radial distributions in units of the Bohr radius.
Probability per radial interval includes the Jacobian $r^2$.}
\end{figure}}{}
\IfFileExists{figures/spherical_box.pdf}{
\begin{figure}[htbp]\centering
\includegraphics[width=.86\linewidth]{figures/spherical_box.pdf}
\caption{Zeros of spherical Bessel functions $j_\ell$ and energy
quantization for a spherical hard-wall box (Problem 4.10).}
\end{figure}}{}
\clearpage
\section{Reference and scope}
D. J. Griffiths and D. F. Schroeter,
\emph{Introduction to Quantum Mechanics}, 3rd ed.,
Cambridge University Press (2018), Chapter 4,
DOI \href{https://doi.org/10.1017/9781316995433}{10.1017/9781316995433}.
These are original educational explanations, not an official solutions
manual; the numbering follows the third-edition textbook.
\end{document}